\scanpage{31}
11.9.1967

Dear Anantharaman,

Matsumura proved that if $X$ is proper over a field $k$, then $\underline{\operatorname{Aut}}_{X/k}$ is representable by a group-scheme locally of finite type over $k$. I think I can systematize the key step of his argument in the following way. Consider a scheme $S$, and a morphism
\[
\phi:Z\longrightarrow X
\]
of $S$-schémes which are proper, flat and of finite presentation. Let $Y$ be locally of finite présentation and separated over $S$, then $\phi$ induces a homomorphisme of functors $u\leadsto u\phi$:
\[
\phi':\underline{\operatorname{Hom}}_S(X,Y)\longrightarrow\underline{\operatorname{Hom}}_S(Z,Y).
\]
Then one can défin\unsure{e} a subfunctor of $\underline{\operatorname{Hom}}_S(\unsure{X},Y)$ where $\phi'$ is ``unramified'' in a rather obvious sense, and this turns out to be an ``open subfunctor'', say $\underline{\operatorname{Hom}}_S(X,Y;\phi)$. Now look at the induced homomorphism
\[
\underline{\operatorname{Hom}}_S(X,Y;\phi)\longrightarrow\underline{\operatorname{Hom}}_S(Z,Y).
\]
Using the main result of Murre's talk, one can prove that the latter morphism is representable by unramified separated morphisms locally of finite presentation; as a consequence, if $\underline{\operatorname{Hom}}_S(Z,Y)$ is representable, so is $\underline{\operatorname{Hom}}_S(X,Y;\phi)$.

To get, given $X$ and $Y$, a representability theorem for $\underline{\operatorname{Hom}}_S(X,Y)$, \unsure{one|onem} tries to find morphisms $\phi_i:Z_i\longrightarrow X$ as above, such that the open subfunctors $\underline{\operatorname{Hom}}_S(X,Y;\phi_i)$ cover $\underline{\operatorname{Hom}}_S(X,Y)$ (as a fpqc sheaf), and such that the \unsure{functorz|functors} $\underline{\operatorname{Hom}}_S(\unsure{Z_i},Y)$ are all representable. If for instance $S$ is the spectrum of a field $k$, and if $X$ has ``enough'' points radicial over $k$ (which is always true if $k$ is alg. closed) then we can take for $Z_i$ all finite \unsure{subshhemes|subschemes} of $X$ whose points are radicial over $k$, and we get
% archive foliation: 30
